\section{逐题单点推导手册}

这一节专门解决“每一题思路都要讲清楚”的问题。它不重复完整代码，而是要求每道题都从自己的题面出发，说明暴力基线、优化线索、状态不变量、边界和变体。复习时可以把这一节当作口述题解提纲。

\subsubsection{LeetCode 1 两数之和}

\textbf{题面模型。} 当前元素只需要寻找唯一补数，历史值到下标的映射就是最小状态。

\textbf{暴力瓶颈。} 先查再插入，避免同一个下标被用两次；若数组有序才考虑双指针。

\textbf{边界反例。} 重复数字、目标为偶数且补数等于当前值、没有答案时返回空结果。

\textbf{变体迁移。} 若改成返回所有不重复数对，要先排序或用频次表处理重复。

\subsubsection{LeetCode 2 两数相加}

\textbf{题面模型。} 链表按低位到高位保存数字，本质是竖式加法而不是整数转换。

\textbf{暴力瓶颈。} 维护两个游标和进位，任一链未结束或进位非零都继续生成节点。

\textbf{边界反例。} 两链长度不同、最后进位、输入为零、不能用内置整数承载超长数字。

\textbf{变体迁移。} 若链表高位在前，可以用栈反向处理或先反转链表。

\subsubsection{LeetCode 3 无重复字符的最长子串}

\textbf{题面模型。} 窗口保存当前无重复子串，右端每次加入一个字符。

\textbf{暴力瓶颈。} 用上次出现位置直接跳过无效左端，左边界只能向右不能回退。

\textbf{边界反例。} 空串、全相同字符、重复字符出现在窗口左侧之外、扩展字符集。

\textbf{变体迁移。} 若要求恰好包含 k 种字符，窗口条件从无重复变成种类计数。

\subsubsection{LeetCode 5 最长回文子串}

\textbf{题面模型。} 回文围绕中心对称，中心可以是字符也可以是两个字符之间。

\textbf{暴力瓶颈。} 中心扩展避免枚举子串后再逐个检查；每次扩展只比较新增左右边界。

\textbf{边界反例。} 偶数长度回文、单字符、全相同字符、多个同长答案。

\textbf{变体迁移。} 若要求回文子序列，应转成区间 DP，不能用中心扩展。

\subsubsection{LeetCode 11 盛最多水的容器}

\textbf{题面模型。} 左右指针代表当前选取的两条边，面积由短板和宽度共同决定。

\textbf{暴力瓶颈。} 每次移动短板，因为移动长板只会缩小宽度且不会提高短板。

\textbf{边界反例。} 两端高度相等、数组长度为二、面积乘法可能溢出。

\textbf{变体迁移。} 接雨水计算每个位置贡献，不能把本题双指针证明直接搬过去。

\subsubsection{LeetCode 15 三数之和}

\textbf{题面模型。} 排序后固定第一个数，剩余两数转成有序双指针。

\textbf{暴力瓶颈。} 和太小移动左指针，和太大移动右指针；固定值和左右值都要跳过重复。

\textbf{边界反例。} 全零、重复元素很多、没有答案、结果顺序不要求但组合不能重复。

\textbf{变体迁移。} 四数之和再外层固定一个数，并用 long long 防止目标和溢出。

\subsubsection{LeetCode 16 最接近的三数之和}

\textbf{题面模型。} 排序后固定一数，双指针寻找最接近目标的两数和。

\textbf{暴力瓶颈。} 每次根据当前和与目标的大小移动指针，同时更新最小绝对差。

\textbf{边界反例。} 差值相等、负数、目标很大、三数和溢出。

\textbf{变体迁移。} 若要求返回所有最接近组合，需要记录同差值并去重。

\subsubsection{LeetCode 18 四数之和}

\textbf{题面模型。} 两层固定加一层双指针，关键是剪枝和去重。

\textbf{暴力瓶颈。} 排序提供单调性，外层可用最小可能和和最大可能和提前跳过。

\textbf{边界反例。} 重复元素、目标超出 int、答案很多导致输出空间主导复杂度。

\textbf{变体迁移。} k 数之和可以递归降维，但工程实现要控制重复和边界。

\subsubsection{LeetCode 26 删除有序数组重复项}

\textbf{题面模型。} 有序数组让重复元素相邻，慢指针表示下一个唯一值写入位置。

\textbf{暴力瓶颈。} 快指针扫描，只有发现新值才写入慢指针并推进。

\textbf{边界反例。} 空数组、全重复、无重复、返回长度不等于物理容量。

\textbf{变体迁移。} 若允许每个元素最多保留两次，只需改变写入条件为和前两个比较。

\subsubsection{LeetCode 27 移除元素}

\textbf{题面模型。} 原地过滤，慢指针左侧始终是不等于目标值的稳定前缀。

\textbf{暴力瓶颈。} 保持顺序用快慢指针；不保持顺序可用左右交换减少写入。

\textbf{边界反例。} 所有元素都被移除、没有元素被移除、目标值在尾部密集。

\textbf{变体迁移。} 工程里要先确认是否要求稳定顺序，不同要求对应不同写法。

\subsubsection{LeetCode 31 下一个排列}

\textbf{题面模型。} 字典序结构由最长非递增后缀决定，枢轴是后缀前一个位置。

\textbf{暴力瓶颈。} 从右找第一个上升点，再从右找刚好更大的元素交换，最后反转后缀。

\textbf{边界反例。} 完全降序、重复数字、只有一个元素、交换后后缀必须最小化。

\textbf{变体迁移。} 若要求上一个排列，比较方向和后缀处理对称变化。

\subsubsection{LeetCode 33 搜索旋转排序数组}

\textbf{题面模型。} 旋转数组每次二分至少有一半保持有序。

\textbf{暴力瓶颈。} 判断哪一半有序，再看目标是否落在有序半边范围内。

\textbf{边界反例。} 未旋转、长度一、目标不存在、边界等号、存在重复元素时退化。

\textbf{变体迁移。} 若重复元素很多，需要在无法判断有序半边时收缩边界。

\subsubsection{LeetCode 34 排序数组中查找首尾位置}

\textbf{题面模型。} 首尾位置可以拆成两个边界：第一个大于等于和第一个大于。

\textbf{暴力瓶颈。} 不要找到一个目标后向两边线性扩展，否则重复元素很多时退化。

\textbf{边界反例。} 目标不存在、目标在开头或结尾、数组为空、全是目标。

\textbf{变体迁移。} 这个模型能迁移到计数某值出现次数和插入区间边界。

\subsubsection{LeetCode 35 搜索插入位置}

\textbf{题面模型。} 题目等价于寻找第一个大于等于目标的位置。

\textbf{暴力瓶颈。} 统一用 lower bound 语义，不在相等时提前返回也能保持边界一致。

\textbf{边界反例。} 目标小于所有元素、目标大于所有元素、重复值、空数组。

\textbf{变体迁移。} 手写二分时用左闭右开可直接返回 left。

\subsubsection{LeetCode 42 接雨水}

\textbf{题面模型。} 每个位置的水由左侧最高和右侧最高的较小者决定。

\textbf{暴力瓶颈。} 前后缀、双指针、单调栈都是不同状态维护方式；要能说明各自结算对象。

\textbf{边界反例。} 单调递增、单调递减、平台高度、多个凹槽、数组长度不足三。

\textbf{变体迁移。} 若柱子变成二维网格，需要最小堆从边界向内扩展。

\subsubsection{LeetCode 49 字母异位词分组}

\textbf{题面模型。} 异位词共享同一个规范表示，可以是排序字符串或字符计数。

\textbf{暴力瓶颈。} 哈希表按规范键分桶，避免两两比较字符串。

\textbf{边界反例。} 空字符串、重复单词、字符集不止小写字母、计数键必须无歧义。

\textbf{变体迁移。} 若字符串很长且字符集固定，计数键通常比排序键更稳定。

\subsubsection{LeetCode 53 最大子数组和}

\textbf{题面模型。} 状态是必须以当前位置结尾的最大连续和。

\textbf{暴力瓶颈。} 若前一个后缀为负，接上它只会拖累当前元素，因此可以重新开始。

\textbf{边界反例。} 全负数组、单元素、和溢出、答案不能初始化为零。

\textbf{变体迁移。} 若要求返回区间下标，同步记录重新开始位置和最优左右端。

\subsubsection{LeetCode 54 螺旋矩阵}

\textbf{题面模型。} 四个边界收缩模拟一圈一圈访问。

\textbf{暴力瓶颈。} 每走完一条边就收缩对应边界，并在走反方向前检查是否仍合法。

\textbf{边界反例。} 单行、单列、空矩阵、最后一圈只剩一个元素。

\textbf{变体迁移。} 若改成生成矩阵，边界逻辑相同，只是读变成写。

\subsubsection{LeetCode 55 跳跃游戏}

\textbf{题面模型。} 扫描过程中只关心当前能到达的最远位置。

\textbf{暴力瓶颈。} 如果下标超过最远可达，任何之前路径都不能到达它；否则用它扩展边界。

\textbf{边界反例。} 第一个元素为零、数组长度一、恰好到达最后、远超最后。

\textbf{变体迁移。} 求最少跳数时把可达区间看成 BFS 层。

\subsubsection{LeetCode 56 合并区间}

\textbf{题面模型。} 排序后可能与当前区间相交的只会是结果尾部。

\textbf{暴力瓶颈。} 按起点排序，扫描时维护结果最后区间的右端点。

\textbf{边界反例。} 端点相接、完全覆盖、无重叠、空列表、输入未排序。

\textbf{变体迁移。} 若区间是半开区间，端点相等不一定合并。

\subsubsection{LeetCode 57 插入区间}

\textbf{题面模型。} 新区间只会影响与它相交的一段连续区间。

\textbf{暴力瓶颈。} 先追加左侧完全不相交区间，再合并重叠段，最后追加右侧。

\textbf{边界反例。} 新区间在最前、最后、覆盖全部、刚好相邻、原列表为空。

\textbf{变体迁移。} 若输入不是有序且不重叠，必须先排序或换模型。

\subsubsection{LeetCode 62 不同路径}

\textbf{题面模型。} 网格状态由上方和左方两种最后一步转移而来。

\textbf{暴力瓶颈。} 二维表可压缩为一维，因为当前行只依赖上一行同列和当前行左侧。

\textbf{边界反例。} 一行、一列、较大网格结果溢出、障碍物版本边界。

\textbf{变体迁移。} 带障碍物时遇到障碍状态清零，第一行第一列要按阻断处理。

\subsubsection{LeetCode 63 不同路径 II}

\textbf{题面模型。} 障碍物让某些格子的路径数变为零。

\textbf{暴力瓶颈。} 滚动数组中遇到障碍直接把当前位置清零，否则累加左侧。

\textbf{边界反例。} 起点或终点是障碍、整行被截断、单行障碍。

\textbf{变体迁移。} 若允许向上或向左走，图会有环，普通填表不再成立。

\subsubsection{LeetCode 64 最小路径和}

\textbf{题面模型。} 状态是到达当前格子的最小代价，最后一步来自上方或左方。

\textbf{暴力瓶颈。} 先初始化第一行和第一列，再处理内部，主转移保持简单。

\textbf{边界反例。} 单行、单列、权值为零、路径和溢出。

\textbf{变体迁移。} 若允许四方向移动，变成最短路而不是普通网格 DP。

\subsubsection{LeetCode 70 爬楼梯}

\textbf{题面模型。} 最后一步来自前一阶或前两阶，是最小的重叠子问题模型。

\textbf{暴力瓶颈。} 滚动变量保存最近两个状态即可，不需要完整数组。

\textbf{边界反例。} n 为一、n 为二、题目是否允许零阶、结果范围。

\textbf{变体迁移。} 若每次可走一到 k 阶，转移变成固定宽度窗口求和。

\subsubsection{LeetCode 72 编辑距离}

\textbf{题面模型。} 二维状态表示两个前缀之间的最少编辑操作。

\textbf{暴力瓶颈。} 末尾字符相同继承左上；不同则枚举插入、删除、替换。

\textbf{边界反例。} 空串、完全相同、完全不同、操作代价不同。

\textbf{变体迁移。} 若只允许插入删除，问题退化到和最长公共子序列相关。

\subsubsection{LeetCode 75 颜色分类}

\textbf{题面模型。} 三指针维护零区、一号区、未知区和二区。

\textbf{暴力瓶颈。} 遇到二时只移动右边界，不移动扫描指针，因为换回来的元素未知。

\textbf{边界反例。} 全零、全二、已经有序、逆序、交换后漏检查。

\textbf{变体迁移。} 这是荷兰国旗问题，能推广到三类分区的原地维护。

\subsubsection{LeetCode 76 最小覆盖子串}

\textbf{题面模型。} 窗口内字符计数必须覆盖目标计数，满足后尽量收缩左端。

\textbf{暴力瓶颈。} 用 formed 记录满足需求的字符种类，避免每次全量比较。

\textbf{边界反例。} 目标有重复字符、源串缺字符、最小窗口在开头或结尾、大小写。

\textbf{变体迁移。} 若字符集固定，数组计数更快；若是 Unicode，哈希表更通用。

\subsubsection{LeetCode 78 子集}

\textbf{题面模型。} 每个元素都有选或不选两种选择，答案是所有路径节点。

\textbf{暴力瓶颈。} 回溯按起点向后枚举，避免不同顺序生成同一集合。

\textbf{边界反例。} 空数组、单元素、输出数量为二的 n 次方、顺序不要求。

\textbf{变体迁移。} 也可用位掩码枚举所有子集，适合 n 较小的场景。

\subsubsection{LeetCode 79 单词搜索}

\textbf{题面模型。} 网格路径不能重复使用同一格，状态包含位置、匹配下标和访问标记。

\textbf{暴力瓶颈。} 剪枝包括首字符不匹配、剩余长度不足、字符频次不足和越界访问。

\textbf{边界反例。} 单字符单词、路径需要回退、重复字母、单元格不能复用。

\textbf{变体迁移。} 若要查很多单词，应使用 Trie 合并搜索前缀。

\subsubsection{LeetCode 84 柱状图中最大的矩形}

\textbf{题面模型。} 每根柱子作为最矮高度时，需要知道左右第一个更矮位置。

\textbf{暴力瓶颈。} 单调递增栈在遇到更矮柱时结算被弹出柱子的最大宽度。

\textbf{边界反例。} 递增、递减、相等高度、哨兵零、宽度计算。

\textbf{变体迁移。} 最大矩形可以作为二维矩阵最大矩形的行压缩子问题。

\subsubsection{LeetCode 88 合并两个有序数组}

\textbf{题面模型。} 第一个数组尾部有空位，适合从后往前写入最终最大值。

\textbf{暴力瓶颈。} 逆向填充避免覆盖尚未读取的 nums1 元素。

\textbf{边界反例。} 第二个数组为空、第一个有效元素为空、重复值、尾部剩余。

\textbf{变体迁移。} 若没有额外空位，只能新开数组或使用更复杂原地归并。

\subsubsection{LeetCode 90 子集 II}

\textbf{题面模型。} 重复元素会让不同下标路径生成相同集合。

\textbf{暴力瓶颈。} 排序后在同一层跳过相同值，保留不同层使用重复值的可能。

\textbf{边界反例。} 全重复、空集、重复值相邻、去重条件写成同层而不是全局。

\textbf{变体迁移。} 组合总和 II 使用同样的同层去重思想。

\subsubsection{LeetCode 94 二叉树中序遍历}

\textbf{题面模型。} 中序访问顺序是左子树、根、右子树，BST 中会得到有序序列。

\textbf{暴力瓶颈。} 用显式栈模拟一路向左，再回退访问根并转向右子树。

\textbf{边界反例。} 空树、单节点、链式左树、链式右树。

\textbf{变体迁移。} Morris 遍历可做到常数额外空间，但会临时改树结构。

\subsubsection{LeetCode 98 验证二叉搜索树}

\textbf{题面模型。} BST 约束是全局上下界，不是只比较父子节点。

\textbf{暴力瓶颈。} 中序严格递增或显式传递上下界都可以；中序版本要保存前一个访问值。

\textbf{边界反例。} 重复值、int 最小最大值、局部合法但全局非法。

\textbf{变体迁移。} 若允许重复值，需要明确重复放左还是放右。

\subsubsection{LeetCode 102 二叉树层序遍历}

\textbf{题面模型。} 队列在每轮开始时保存当前层所有节点。

\textbf{暴力瓶颈。} 记录 level size 固定层边界，避免下一层节点混入当前层。

\textbf{边界反例。} 空树、只有根、最后一层不满、左右顺序。

\textbf{变体迁移。} 锯齿形层序只是在每层输出顺序上增加方向控制。

\subsubsection{LeetCode 103 二叉树锯齿形层序遍历}

\textbf{题面模型。} BFS 层边界不变，变化的是当前层写入结果的方向。

\textbf{暴力瓶颈。} 可以层内反转，也可以预分配数组按方向写入目标下标。

\textbf{边界反例。} 单层、偶数层、奇数层、空树。

\textbf{变体迁移。} 若要求从底向上输出，收集后反转层列表即可。

\subsubsection{LeetCode 104 二叉树最大深度}

\textbf{题面模型。} 深度可以看成层数，也可以作为栈中携带的状态。

\textbf{暴力瓶颈。} 显式栈保存节点和深度，避免极端链式树递归栈过深。

\textbf{边界反例。} 空树返回零、根深度是一、链式树。

\textbf{变体迁移。} 最小深度不能简单取左右最小，因为空子树不构成叶子路径。

\subsubsection{LeetCode 105 前序中序构造二叉树}

\textbf{题面模型。} 前序给根，中序把左右子树切开。

\textbf{暴力瓶颈。} 用哈希表记录中序下标，避免每次线性寻找根位置。

\textbf{边界反例。} 节点值唯一、空子树、左右子树长度计算、输入不合法。

\textbf{变体迁移。} 后序中序构造类似，只是根来自后序末尾。

\subsubsection{LeetCode 106 中序后序构造二叉树}

\textbf{题面模型。} 后序末尾是根，中序位置确定左右子树规模。

\textbf{暴力瓶颈。} 构造时要小心后序左右区间边界，避免切分错位。

\textbf{边界反例。} 空树、单节点、全左链、全右链。

\textbf{变体迁移。} 若给前序和后序，二叉树通常不能唯一确定，除非额外限制。

\subsubsection{LeetCode 108 有序数组转二叉搜索树}

\textbf{题面模型。} 有序数组中点作为根能让左右规模接近。

\textbf{暴力瓶颈。} 每次选择区间中点，左右子区间构造左右子树。

\textbf{边界反例。} 空区间、偶数长度选左中还是右中、高度平衡定义。

\textbf{变体迁移。} 工程上递归可读，若严格避免递归可用队列保存区间和父指针。

\subsubsection{LeetCode 110 平衡二叉树}

\textbf{题面模型。} 每个节点需要知道子树高度以及是否已经失衡。

\textbf{暴力瓶颈。} 后序汇总时一旦子树失衡就向上返回失败标记，避免重复算高度。

\textbf{边界反例。} 空树、单节点、左右高度差正好一、链式树。

\textbf{变体迁移。} 可以把返回值设计成 optional 高度，空表示不平衡。

\subsubsection{LeetCode 111 二叉树最小深度}

\textbf{题面模型。} 最小深度是到最近叶子的节点数，不是左右深度简单取最小。

\textbf{暴力瓶颈。} BFS 第一次遇到叶子就是最小深度，因为按层扩展。

\textbf{边界反例。} 只有左子树、只有右子树、根为空、根是叶子。

\textbf{变体迁移。} DFS 写法必须特殊处理空子树不能作为最小路径。

\subsubsection{LeetCode 112 路径总和}

\textbf{题面模型。} 路径必须从根到叶子，状态是剩余目标值。

\textbf{暴力瓶颈。} 沿路径减去当前节点值，到叶子时检查剩余是否等于叶子值。

\textbf{边界反例。} 负数节点、目标为零、非叶子提前达到目标、空树。

\textbf{变体迁移。} 若要求列出所有路径，需要保存路径数组并在回退时弹出。

\subsubsection{LeetCode 113 路径总和 II}

\textbf{题面模型。} 相比判定题，状态还要保存当前路径。

\textbf{暴力瓶颈。} 到叶子且和满足时复制路径；回溯时恢复路径。

\textbf{边界反例。} 多个答案、负数、空树、路径数组复制时机。

\textbf{变体迁移。} 若路径不要求从根到叶，模型会变成前缀和加哈希。

\subsubsection{LeetCode 114 二叉树展开为链表}

\textbf{题面模型。} 目标是按前序顺序把树原地改成右指针链。

\textbf{暴力瓶颈。} 显式栈按右后左先入顺序处理，逐个接到前驱右侧。

\textbf{边界反例。} 空树、单节点、原有右子树不能丢、左指针置空。

\textbf{变体迁移。} 也可用寻找左子树最右节点的原地方法。

\subsubsection{LeetCode 121 买卖股票的最佳时机}

\textbf{题面模型。} 卖出日固定时，最佳买入日是此前最低价格。

\textbf{暴力瓶颈。} 扫描维护历史最低价和最大利润，一次交易无需复杂状态机。

\textbf{边界反例。} 单日价格、一直下跌、一直上涨、价格相同。

\textbf{变体迁移。} 多次交易、冷冻期和手续费版本都要扩展成状态机。

\subsubsection{LeetCode 122 买卖股票 II}

\textbf{题面模型。} 允许多次交易且不能同时持有，所有正收益差都可以收集。

\textbf{暴力瓶颈。} 把长上涨段拆成每天买卖收益相同，因此累加正差值最优。

\textbf{边界反例。} 下降段、平台价格、只有一天、手续费不存在。

\textbf{变体迁移。} 有手续费时不能简单累加正差，需要状态机或调整贪心阈值。

\subsubsection{LeetCode 123 买卖股票 III}

\textbf{题面模型。} 最多两次交易，状态需要包含交易次数和持有状态。

\textbf{暴力瓶颈。} 维护 buy1、sell1、buy2、sell2 四个状态，按天更新。

\textbf{边界反例。} 价格为空、只够一次交易、两次交易间不能重叠。

\textbf{变体迁移。} 最多 k 次交易是它的泛化，k 很大时退化为无限次交易。

\subsubsection{LeetCode 124 二叉树最大路径和}

\textbf{题面模型。} 路径可以从任意节点到任意节点，但不能分叉向父节点贡献两边。

\textbf{暴力瓶颈。} 每个节点向父亲贡献单边最大增益，全局答案可用左右两边加当前节点更新。

\textbf{边界反例。} 全负数、单节点、左右增益为负时取零、答案初始值。

\textbf{变体迁移。} 若路径必须从根到叶，状态和答案位置完全不同。

\subsubsection{LeetCode 127 单词接龙}

\textbf{题面模型。} 每个单词是节点，一次修改一个字符形成边。

\textbf{暴力瓶颈。} 通配模式桶加速邻居查询，BFS 第一次到达终点就是最短转换长度。

\textbf{边界反例。} 终点不在词表、起点和终点相同、桶重复扫描、单词长度一致。

\textbf{变体迁移。} 双向 BFS 可以进一步降低搜索空间。

\subsubsection{LeetCode 128 最长连续序列}

\textbf{题面模型。} 连续段只需要从没有前驱的数开始扩展。

\textbf{暴力瓶颈。} 哈希集合判断 x 减一 是否存在，跳过非起点避免重复扫描。

\textbf{边界反例。} 重复元素、负数、空数组、整数边界。

\textbf{变体迁移。} 若数据流动态插入，可以用并查集或区间合并维护长度。

\subsubsection{LeetCode 130 被围绕的区域}

\textbf{题面模型。} 只有不连通到边界的 O 才会被翻转。

\textbf{暴力瓶颈。} 从边界 O 出发标记安全区域，再翻转剩余 O。

\textbf{边界反例。} 全边界、无 O、单行单列、内部岛屿。

\textbf{变体迁移。} 这是反向思维：找不能被翻转的区域比直接判断每块区域更简单。

\subsubsection{LeetCode 131 分割回文串}

\textbf{题面模型。} 每次选择一个从当前位置开始的回文前缀。

\textbf{暴力瓶颈。} 预处理回文 DP 表，让回溯合法性检查为常数。

\textbf{边界反例。} 空串、单字符、全相同字符、重复分割路径。

\textbf{变体迁移。} 若只求最少切割次数，可改成 DP 最短切分。

\subsubsection{LeetCode 133 克隆图}

\textbf{题面模型。} 每个原节点需要唯一对应一个新节点，边按原图连接。

\textbf{暴力瓶颈。} 哈希表保存原节点到克隆节点的映射，BFS 或 DFS 遍历图。

\textbf{边界反例。} 空图、自环、重边、环结构、不能按节点值唯一假设。

\textbf{变体迁移。} 若图节点值不唯一，必须用地址作为键。

\subsubsection{LeetCode 136 只出现一次的数字}

\textbf{题面模型。} 成对元素可以用异或抵消，剩下单个元素。

\textbf{暴力瓶颈。} 异或满足交换结合，扫描顺序不影响答案。

\textbf{边界反例。} 零、负数、唯一值在开头或结尾、所有其他值恰好两次。

\textbf{变体迁移。} 若其他值出现三次，要改为按位计数取模。

\subsubsection{LeetCode 139 单词拆分}

\textbf{题面模型。} 状态表示前缀能否由字典单词拼出。

\textbf{暴力瓶颈。} 枚举最后一个单词的起点，左侧前缀可达且子串在字典中则当前可达。

\textbf{边界反例。} 空串、重复单词、长单词、substr 复制成本。

\textbf{变体迁移。} 若要输出所有句子，需要 DFS 加记忆化或前驱列表。

\subsubsection{LeetCode 141 环形链表}

\textbf{题面模型。} 快慢指针在有环时必然相遇，无环时快指针先到空。

\textbf{暴力瓶颈。} 不用哈希表也能常数空间检测环。

\textbf{边界反例。} 空链、单节点自环、两节点环、尾部无环。

\textbf{变体迁移。} 若要求环入口，用相遇后双指针同步走。

\subsubsection{LeetCode 142 环形链表 II}

\textbf{题面模型。} 相遇点到入口的距离关系来自快慢指针速度差。

\textbf{暴力瓶颈。} 相遇后一个指针回头，两个指针每次走一步，再次相遇就是入口。

\textbf{边界反例。} 入口在头节点、长尾短环、无环、单节点环。

\textbf{变体迁移。} 也可用哈希集合，但空间更高。

\subsubsection{LeetCode 146 LRU 缓存}

\textbf{题面模型。} 哈希表负责按 key 定位，双向链表负责维护新旧顺序。

\textbf{暴力瓶颈。} 访问和更新都把节点移动到头部，容量满时删除尾部最旧节点。

\textbf{边界反例。} 容量为一、重复 put、get 不存在、更新已有值。

\textbf{变体迁移。} C++ 可用 \code{std::list} 加迭代器，也可手写双向链表理解所有权。

\subsubsection{LeetCode 148 排序链表}

\textbf{题面模型。} 链表不能随机访问，适合归并排序而不是快速排序数组 partition。

\textbf{暴力瓶颈。} 自底向上迭代归并能避免递归栈，同时保持 O(n log n)。

\textbf{边界反例。} 空链、单节点、重复值、奇数长度、切断子链。

\textbf{变体迁移。} 若把值复制到数组排序，会改变节点身份语义。

\subsubsection{LeetCode 150 逆波兰表达式求值}

\textbf{题面模型。} 后缀表达式把优先级和括号都编码进 token 顺序。

\textbf{暴力瓶颈。} 遇到数字入栈，遇到运算符弹出右操作数和左操作数计算。

\textbf{边界反例。} 负数 token、多位数、除法向零截断、减法顺序。

\textbf{变体迁移。} 中缀表达式求值需要另一个运算符栈处理优先级。

\subsubsection{LeetCode 152 乘积最大子数组}

\textbf{题面模型。} 负数会让最大乘积和最小乘积互换。

\textbf{暴力瓶颈。} 同时维护以当前位置结尾的最大和最小乘积。

\textbf{边界反例。} 零、负数个数奇偶、单元素、乘积溢出。

\textbf{变体迁移。} 如果要求返回区间，同步记录状态来源。

\subsubsection{LeetCode 153 寻找旋转排序数组中的最小值}

\textbf{题面模型。} 最小值是旋转断点，右端点可帮助判断中点在哪一段。

\textbf{暴力瓶颈。} 若 mid 大于 right，最小值在右侧；否则在左侧含 mid。

\textbf{边界反例。} 未旋转、长度一、旋转一位、无重复条件。

\textbf{变体迁移。} 有重复时遇到相等只能收缩 right，最坏退化。

\subsubsection{LeetCode 154 寻找旋转排序数组中的最小值 II}

\textbf{题面模型。} 重复元素会让 mid 和 right 相等时无法判断最小值方向。

\textbf{暴力瓶颈。} 相等时只能安全地丢掉一个 right，因为它不是唯一必要候选。

\textbf{边界反例。} 全相等、重复值跨断点、最小值出现多次。

\textbf{变体迁移。} 这题说明二分依赖的信息不足时会退化，不是所有旋转题都严格对数。

\subsubsection{LeetCode 155 最小栈}

\textbf{题面模型。} 普通栈只能看到栈顶，不能常数时间知道历史最小值。

\textbf{暴力瓶颈。} 辅助栈同步保存每个深度对应的最小值。

\textbf{边界反例。} 重复最小值、弹出最小值、空栈操作约定。

\textbf{变体迁移。} 也可保存差值压缩空间，但可读性和溢出处理更复杂。

\subsubsection{LeetCode 160 相交链表}

\textbf{题面模型。} 相交是节点地址相同，不是节点值相同。

\textbf{暴力瓶颈。} 两个指针走完各自链表后切换到另一条链，走过总长度相同后相遇。

\textbf{边界反例。} 不相交、头节点相交、尾节点相交、长度差很大。

\textbf{变体迁移。} 哈希集合更直观但空间更高。

\subsubsection{LeetCode 162 寻找峰值}

\textbf{题面模型。} 比较 mid 和 mid 加一 的斜率能判断某侧一定存在峰值。

\textbf{暴力瓶颈。} 上升则右侧有峰，下降则左侧含 mid 有峰。

\textbf{边界反例。} 长度一、峰在边界、多个峰、相邻不相等假设。

\textbf{变体迁移。} 二维峰值需要按行列最大值进行更复杂二分。

\subsubsection{LeetCode 169 多数元素}

\textbf{题面模型。} 超过一半的元素可以和其他元素两两抵消后仍然剩下。

\textbf{暴力瓶颈。} Boyer-Moore 投票维护候选和计数。

\textbf{边界反例。} 题目是否保证存在多数、计数归零、全相同。

\textbf{变体迁移。} 若找超过三分之一的元素，需要维护两个候选。

\subsubsection{LeetCode 189 轮转数组}

\textbf{题面模型。} 右移 k 位等价于把数组切成两段后换序。

\textbf{暴力瓶颈。} 三次反转能原地完成：整体、前 k、后 n 减 k。

\textbf{边界反例。} k 为零、k 大于 n、空数组、单元素。

\textbf{变体迁移。} 环状替换也可做，但要处理最大公约数个循环。

\subsubsection{LeetCode 198 打家劫舍}

\textbf{题面模型。} 每间房只有偷和不偷两种结局，偷当前就不能偷前一间。

\textbf{暴力瓶颈。} 滚动保存前一状态和前两状态。

\textbf{边界反例。} 空数组、单间、全零、金额很大。

\textbf{变体迁移。} 环形房屋要拆成不含首和不含尾两个线性问题。

\subsubsection{LeetCode 199 二叉树右视图}

\textbf{题面模型。} 每层从右侧能看到的就是该层最后访问或最先访问的右侧节点。

\textbf{暴力瓶颈。} BFS 固定层边界，取每层最后一个；或 DFS 按右优先首次到达深度。

\textbf{边界反例。} 空树、左链、右链、左右交错。

\textbf{变体迁移。} 若求左视图，只改变层内取值方向。

\subsubsection{LeetCode 200 岛屿数量}

\textbf{题面模型。} 陆地格子是节点，四方向相邻是边，岛屿是连通块。

\textbf{暴力瓶颈。} 扫描遇到未访问陆地就启动搜索并标记整块。

\textbf{边界反例。} 空网格、全水、全陆、斜对角不连通、是否允许修改输入。

\textbf{变体迁移。} 也可用并查集合并相邻陆地，适合动态岛屿扩展。

\subsubsection{LeetCode 202 快乐数}

\textbf{题面模型。} 反复变换数字会进入一或循环，状态空间最终有限。

\textbf{暴力瓶颈。} 用集合检测重复状态，或用快慢指针看成函数图。

\textbf{边界反例。} 输入为一、进入循环、位平方和函数、负数不在题意内。

\textbf{变体迁移。} 快慢指针版本能训练把数值迭代看成链表。

\subsubsection{LeetCode 206 反转链表}

\textbf{题面模型。} 每次把当前节点的 next 指向已经反转好的前缀头。

\textbf{暴力瓶颈。} 先保存后继，再改指针，再推进 previous 和 current。

\textbf{边界反例。} 空链、单节点、两个节点、不要丢失剩余链。

\textbf{变体迁移。} K 组反转是在这段局部反转外增加分组边界。

\subsubsection{LeetCode 207 课程表}

\textbf{题面模型。} 课程是节点，先修关系是有向边，问题是判断是否有环。

\textbf{暴力瓶颈。} 入度拓扑排序不断学习当前无依赖课程。

\textbf{边界反例。} 边方向、孤立课程、环、自依赖、重复边。

\textbf{变体迁移。} DFS 三色标记也能判环，但要明确访问状态。

\subsubsection{LeetCode 208 实现 Trie}

\textbf{题面模型。} 前缀树把字符串公共前缀共享成路径。

\textbf{暴力瓶颈。} 插入和查询都逐字符沿边移动，不存在的边按需要创建。

\textbf{边界反例。} 空字符串约定、单词结束标记、前缀存在不等于单词存在。

\textbf{变体迁移。} 若字符集很大，用哈希子节点；小写字母用数组更快。

\subsubsection{LeetCode 209 长度最小的子数组}

\textbf{题面模型。} 正数数组让窗口和随右端增加、随左端减少。

\textbf{暴力瓶颈。} 右端扩展到满足目标后，不断左移缩短并更新答案。

\textbf{边界反例。} 不存在答案、单元素达到目标、全正是必要条件。

\textbf{变体迁移。} 若含负数，需要前缀和加单调队列。

\subsubsection{LeetCode 210 课程表 II}

\textbf{题面模型。} 在判环基础上还要输出一个可行拓扑序。

\textbf{暴力瓶颈。} 每弹出一个入度为零课程，就把它加入顺序并删除出边。

\textbf{边界反例。} 有环返回空数组、多个可行顺序、边方向。

\textbf{变体迁移。} 若需要字典序最小拓扑序，把队列换成小顶堆。

\subsubsection{LeetCode 211 添加与搜索单词}

\textbf{题面模型。} 点号通配符让查询可能分叉，但前缀树仍能剪掉不匹配前缀。

\textbf{暴力瓶颈。} 普通字符沿唯一边走，点号枚举当前节点所有孩子。

\textbf{边界反例。} 连续通配符、空结果、单词结束标记、字符集大小。

\textbf{变体迁移。} 大量通配符会接近遍历整棵 Trie，需要规模约束。

\subsubsection{LeetCode 215 数组中的第 K 个最大元素}

\textbf{题面模型。} 只关心第 K 大，不需要完整排序。

\textbf{暴力瓶颈。} 大小为 K 的小顶堆保存当前前 K 大边界；快速选择可做到平均线性。

\textbf{边界反例。} K 为一、K 为 n、重复值、随机 pivot。

\textbf{变体迁移。} 若数据流持续到来，堆方案比一次性快速选择更自然。

\subsubsection{LeetCode 216 组合总和 III}

\textbf{题面模型。} 从一到九中选 k 个数和为 n，每个数最多用一次。

\textbf{暴力瓶颈。} 按递增起点枚举，利用剩余个数和剩余和剪枝。

\textbf{边界反例。} n 太小、n 太大、k 为零、路径长度已经超过 k。

\textbf{变体迁移。} 可预估剩余最小和最大和进一步剪枝。

\subsubsection{LeetCode 217 存在重复元素}

\textbf{题面模型。} 判断是否出现过是哈希集合最直接的职责。

\textbf{暴力瓶颈。} 扫描时若插入前已经存在则立即返回真。

\textbf{边界反例。} 空数组、全唯一、重复在末尾、值域很小。

\textbf{变体迁移。} 若允许修改输入，排序后看相邻可以降低额外空间。

\subsubsection{LeetCode 221 最大正方形}

\textbf{题面模型。} 以当前位置为右下角的最大正方形由上、左、左上三个状态限制。

\textbf{暴力瓶颈。} 当前位置为一时取三者最小加一，否则为零。

\textbf{边界反例。} 全零、全一、单行单列、字符一和数字一不要混淆。

\textbf{变体迁移。} 最大矩形则需要按行构造柱状图再用单调栈。

\subsubsection{LeetCode 226 翻转二叉树}

\textbf{题面模型。} 每个节点交换左右孩子即可，遍历顺序不影响最终结构。

\textbf{暴力瓶颈。} 显式队列或栈逐个节点交换，避免递归深度。

\textbf{边界反例。} 空树、单节点、只有一侧子树、重复值。

\textbf{变体迁移。} 这题简单但适合训练树节点指针修改。

\subsubsection{LeetCode 230 二叉搜索树中第 K 小}

\textbf{题面模型。} BST 中序遍历按升序访问节点。

\textbf{暴力瓶颈。} 显式栈中序，访问第 k 个节点时返回。

\textbf{边界反例。} k 为一、k 为节点数、树不平衡、是否有重复值。

\textbf{变体迁移。} 若频繁查询，可在节点维护子树大小成为顺序统计树。

\subsubsection{LeetCode 232 用栈实现队列}

\textbf{题面模型。} 两个栈分别负责输入顺序和输出顺序。

\textbf{暴力瓶颈。} 只有输出栈为空时，才把输入栈整体倒过去，保证均摊常数。

\textbf{边界反例。} 连续 peek、空队列、交替 push pop。

\textbf{变体迁移。} 用队列实现栈则要通过轮转保持最近元素在队首。

\subsubsection{LeetCode 235 二叉搜索树最近公共祖先}

\textbf{题面模型。} BST 的值域关系能决定两个目标在当前节点的哪一侧。

\textbf{暴力瓶颈。} 两个值都小去左，都大去右，否则当前节点就是分叉点。

\textbf{边界反例。} 一个节点是另一个祖先、目标顺序、重复值约定。

\textbf{变体迁移。} 普通二叉树不能用值域剪枝，需要父指针或后序状态。

\subsubsection{LeetCode 236 二叉树最近公共祖先}

\textbf{题面模型。} 普通树没有有序性，只能依靠祖先关系。

\textbf{暴力瓶颈。} 父指针集合或后序返回目标命中状态都可行。

\textbf{边界反例。} 目标不存在、p 等于 q、一个是另一个祖先、比较节点地址。

\textbf{变体迁移。} 多次 LCA 查询可预处理倍增祖先表。

\subsubsection{LeetCode 238 除自身以外数组的乘积}

\textbf{题面模型。} 不能用除法时，答案由左侧乘积和右侧乘积组成。

\textbf{暴力瓶颈。} 先写入左侧前缀积，再从右向左乘右侧后缀积。

\textbf{边界反例。} 一个零、多个零、负数、乘积溢出。

\textbf{变体迁移。} 这个题展示输出数组可作为状态存储的一部分。

\subsubsection{LeetCode 239 滑动窗口最大值}

\textbf{题面模型。} 每个窗口最大值来自仍未过期且没有被新元素支配的候选。

\textbf{暴力瓶颈。} 单调队列保存下标，队尾小于等于新值的候选被弹出。

\textbf{边界反例。} k 为一、k 等于数组长度、重复最大值、队首过期。

\textbf{变体迁移。} 受限子序列和使用同样队列维护最近 k 个 DP 最大值。

\subsubsection{LeetCode 240 搜索二维矩阵 II}

\textbf{题面模型。} 右上角同时具有向左变小、向下变大的排除能力。

\textbf{暴力瓶颈。} 当前值大于目标就左移，小于目标就下移。

\textbf{边界反例。} 空矩阵、单行、单列、目标在角落、重复值。

\textbf{变体迁移。} 从左上角无法一次排除一行或一列。

\subsubsection{LeetCode 279 完全平方数}

\textbf{题面模型。} 完全平方数可以看成无限使用的物品，目标是凑出 n 的最少数量。

\textbf{暴力瓶颈。} 完全背包最少物品数，或把数字看成图节点做 BFS。

\textbf{边界反例。} n 为零或一、完全平方数本身、较大 n。

\textbf{变体迁移。} 数学四平方定理可给更强解法，但面试 DP 更通用。

\subsubsection{LeetCode 283 移动零}

\textbf{题面模型。} 稳定保留非零元素顺序，把零集中到尾部。

\textbf{暴力瓶颈。} 慢指针写入下一个非零位置，最后补零或用交换。

\textbf{边界反例。} 全零、无零、零在开头结尾、稳定顺序要求。

\textbf{变体迁移。} 若目标是移动指定值，模型与移除元素相同。

\subsubsection{LeetCode 287 寻找重复数}

\textbf{题面模型。} 数组值域使下标到值形成函数图，重复数是环入口。

\textbf{暴力瓶颈。} 快慢指针不修改数组且常数空间；值域二分用抽屉原理。

\textbf{边界反例。} 重复数出现多次、不能排序、不能改数组、长度为 n 加一。

\textbf{变体迁移。} 快慢指针要求值都能作为合法下标。

\subsubsection{LeetCode 295 数据流的中位数}

\textbf{题面模型。} 数据流需要动态维护较小一半和较大一半。

\textbf{暴力瓶颈。} 大顶堆保存较小一半，小顶堆保存较大一半，大小差不超过一。

\textbf{边界反例。} 偶数个元素、负数、重复值、平均值溢出。

\textbf{变体迁移。} 若还要删除任意元素，需要延迟删除哈希表。

\subsubsection{LeetCode 300 最长递增子序列}

\textbf{题面模型。} 二次 DP 固定结尾，二分优化维护每个长度的最小结尾。

\textbf{暴力瓶颈。} tails 不是实际答案序列，而是未来扩展潜力表。

\textbf{边界反例。} 重复元素、严格递增和非递减区别、空数组。

\textbf{变体迁移。} 若要还原路径，需要额外前驱数组，不能只看 tails。

\subsubsection{LeetCode 309 最佳买卖股票含冷冻期}

\textbf{题面模型。} 冷冻期让买入来源受昨天卖出状态限制。

\textbf{暴力瓶颈。} 状态机维护持有、刚卖出、休息三种日终状态。

\textbf{边界反例。} 空价格、一天、连续上涨、卖出后不能次日买入。

\textbf{变体迁移。} 手续费版本可以在买入或卖出转移中扣费。

\subsubsection{LeetCode 322 零钱兑换}

\textbf{题面模型。} 金额是容量，硬币可无限使用，目标是最少硬币数。

\textbf{暴力瓶颈。} dp[value] 从 value 减 coin 的已知状态转移。

\textbf{边界反例。} 无法凑出、amount 为零、硬币面额大于目标、哨兵溢出。

\textbf{变体迁移。} 零钱兑换 II 问组合数，循环语义不同。

\subsubsection{LeetCode 337 打家劫舍 III}

\textbf{题面模型。} 树上相邻节点不能同时选，每个节点需要偷和不偷两个状态。

\textbf{暴力瓶颈。} 后序汇总：偷当前则孩子不能偷，不偷当前则孩子可偷可不偷。

\textbf{边界反例。} 空树、负值是否可能、链式树、递归深度。

\textbf{变体迁移。} 可用显式栈后序保存节点状态，避免调用栈。

\subsubsection{LeetCode 347 前 K 个高频元素}

\textbf{题面模型。} 先统计频次，再只维护前 K 个候选。

\textbf{暴力瓶颈。} 小顶堆大小超过 K 就弹出最低频；桶排序利用频次上界。

\textbf{边界反例。} K 等于不同元素数、频次相同、结果顺序不要求。

\textbf{变体迁移。} 若要按频次和字典序排序，比较器要同时处理两个键。

\subsubsection{LeetCode 394 字符串解码}

\textbf{题面模型。} 嵌套结构需要保存上一层字符串和重复次数。

\textbf{暴力瓶颈。} 遇到左括号入栈当前状态，右括号弹出并拼接展开。

\textbf{边界反例。} 多位数字、嵌套、空片段、数字后必须跟括号。

\textbf{变体迁移。} 递归下降解析更通用，但迭代栈更符合工程栈控制。

\subsubsection{LeetCode 416 分割等和子集}

\textbf{题面模型。} 总和一半是背包容量，每个数只能使用一次。

\textbf{暴力瓶颈。} 0-1 背包可达性，一维容量必须倒序。

\textbf{边界反例。} 总和为奇数、目标为零、数值较大、复杂度依赖 target。

\textbf{变体迁移。} 负数会破坏普通容量模型，需要重新建模。

\subsubsection{LeetCode 417 太平洋大西洋水流问题}

\textbf{题面模型。} 从每个格子向海搜索很慢，反向从海边沿非递减高度搜索。

\textbf{暴力瓶颈。} 分别标记能到太平洋和大西洋的格子，交集就是答案。

\textbf{边界反例。} 单行单列、平台高度、边界重复入队、方向条件反向。

\textbf{变体迁移。} 反向搜索是很多可达性题的重要技巧。

\subsubsection{LeetCode 435 无重叠区间}

\textbf{题面模型。} 保留最多不重叠区间等价于删除最少区间。

\textbf{暴力瓶颈。} 按终点升序选择，结束越早给后续留下空间越多。

\textbf{边界反例。} 端点相接是否重叠、完全覆盖、相同终点、空列表。

\textbf{变体迁移。} 用交换论证说明最早结束的选择不劣于其他第一选择。

\subsubsection{LeetCode 438 找到字符串中所有字母异位词}

\textbf{题面模型。} 固定长度窗口内字符频次等于目标频次即为答案。

\textbf{暴力瓶颈。} 右进左出维护计数，窗口长度达到目标长度后检查差异。

\textbf{边界反例。} 目标比源串长、重复字符、字符集固定、答案重叠。

\textbf{变体迁移。} 维护差异计数可以避免每次比较整个数组。

\subsubsection{LeetCode 450 删除二叉搜索树中的节点}

\textbf{题面模型。} 删除节点后仍要保持 BST 中序有序。

\textbf{暴力瓶颈。} 若有两个孩子，用后继或前驱替换当前值，再删除那个后继节点。

\textbf{边界反例。} 删除叶子、一个孩子、两个孩子、删除根节点。

\textbf{变体迁移。} 工程实现要清楚是移动值还是重连节点。

\subsubsection{LeetCode 452 用最少数量的箭引爆气球}

\textbf{题面模型。} 一支箭的位置可以覆盖所有包含该位置的区间。

\textbf{暴力瓶颈。} 按右端点排序，当前箭放在最早结束气球的右端。

\textbf{边界反例。} 端点相同、负坐标、区间包含、闭区间边界。

\textbf{变体迁移。} 它和无重叠区间是同一类最早结束贪心。

\subsubsection{LeetCode 494 目标和}

\textbf{题面模型。} 正负号选择可代数转换为选若干数凑出特定正和。

\textbf{暴力瓶颈。} 若 S 加 target 为奇数或目标绝对值过大，直接无解。

\textbf{边界反例。} 零元素会让方案数翻倍、目标为负、容量为零。

\textbf{变体迁移。} 也可用哈希 DP 保存所有可能和，但背包转换更高效。

\subsubsection{LeetCode 496 下一个更大元素 I}

\textbf{题面模型。} 第二个数组提供每个值右侧第一个更大值的全局关系。

\textbf{暴力瓶颈。} 单调栈预处理值到下一个更大值映射，再回答第一个数组查询。

\textbf{边界反例。} 不存在更大值、递减数组、值唯一假设。

\textbf{变体迁移。} 若值不唯一，应按下标而不是值建映射。

\subsubsection{LeetCode 503 下一个更大元素 II}

\textbf{题面模型。} 环形数组可以通过遍历两倍下标模拟绕回。

\textbf{暴力瓶颈。} 第一遍入栈，第二遍只帮助未解决下标找到答案。

\textbf{边界反例。} 全递减、全相等、长度一、取模下标。

\textbf{变体迁移。} 不要在第二遍重复入栈，否则可能死循环或重复计算。

\subsubsection{LeetCode 518 零钱兑换 II}

\textbf{题面模型。} 硬币无限使用，问组合数而不是最少数量。

\textbf{暴力瓶颈。} 外层遍历硬币、内层金额正序，保证组合按固定硬币顺序统计。

\textbf{边界反例。} amount 为零、无硬币、组合数较大、循环顺序反例。

\textbf{变体迁移。} 若外层金额内层硬币，会统计排列数。

\subsubsection{LeetCode 543 二叉树直径}

\textbf{题面模型。} 直径经过某节点时等于左高加右高，但向父亲只能贡献单边高度。

\textbf{暴力瓶颈。} 后序遍历同时更新全局直径并返回高度。

\textbf{边界反例。} 空树、单节点、链式树、边数还是节点数定义。

\textbf{变体迁移。} 最大路径和是同型题，但贡献值和全局更新有负数处理。

\subsubsection{LeetCode 547 省份数量}

\textbf{题面模型。} 城市连接关系形成无向图，省份就是连通分量。

\textbf{暴力瓶颈。} 并查集合并所有直接连接的城市，最后统计根数量。

\textbf{边界反例。} 邻接矩阵对称、自连接、孤立城市、只扫描上三角。

\textbf{变体迁移。} BFS 也可做；并查集更突出动态合并。

\subsubsection{LeetCode 560 和为 K 的子数组}

\textbf{题面模型。} 当前前缀和减去目标值，就是需要匹配的历史前缀。

\textbf{暴力瓶颈。} 哈希表保存前缀和出现次数，先查询再更新。

\textbf{边界反例。} 负数、目标为零、从下标零开始的区间、重复前缀。

\textbf{变体迁移。} 若要求最长区间，哈希表应保存最早位置而不是次数。

\subsubsection{LeetCode 581 最短无序连续子数组}

\textbf{题面模型。} 要找破坏全局有序性的最左和最右位置。

\textbf{暴力瓶颈。} 从左维护历史最大值确定右边界，从右维护历史最小值确定左边界。

\textbf{边界反例。} 已经有序、逆序、重复值、局部乱序。

\textbf{变体迁移。} 排序比较法更直观但复杂度更高。

\subsubsection{LeetCode 621 任务调度器}

\textbf{题面模型。} 相同任务之间有冷却间隔，核心受最高频任务骨架限制。

\textbf{暴力瓶颈。} 可用大顶堆模拟，也可用最高频公式计算最短时间。

\textbf{边界反例。} 冷却为零、多个最高频、空闲被其他任务填满。

\textbf{变体迁移。} 若任务有不同释放时间或执行时长，就变成更一般的调度问题。

\subsubsection{LeetCode 647 回文子串}

\textbf{题面模型。} 每个回文都有唯一中心或中心间隙。

\textbf{暴力瓶颈。} 对每个中心向两边扩展，扩展成功一次就计数一次。

\textbf{边界反例。} 空串、全相同字符、偶数回文、奇数回文。

\textbf{变体迁移。} Manacher 算法可线性解决，但中心扩展更适合面试基础。

\subsubsection{LeetCode 684 冗余连接}

\textbf{题面模型。} 树加一条边会形成唯一环。

\textbf{暴力瓶颈。} 按顺序加边，若一条边两端已连通，则它就是冗余边。

\textbf{边界反例。} 节点编号、最后一条成环边、重复边、并查集初始化大小。

\textbf{变体迁移。} 有向冗余连接需要同时处理入度为二和环，复杂度更高。

\subsubsection{LeetCode 695 岛屿的最大面积}

\textbf{题面模型。} 面积就是一个连通块中陆地格子的数量。

\textbf{暴力瓶颈。} 每发现新陆地就搜索完整连通块并计数，取最大值。

\textbf{边界反例。} 全水、全陆、细长岛、是否允许修改输入。

\textbf{变体迁移。} 与岛屿数量相比，只是连通块聚合值不同。

\subsubsection{LeetCode 704 二分查找}

\textbf{题面模型。} 基础题的重点是固定区间语义。

\textbf{暴力瓶颈。} 左闭右开写法能统一 lower bound 和存在性检查。

\textbf{边界反例。} 空数组、长度一、目标在边界、目标不存在。

\textbf{变体迁移。} 所有复杂二分都应该能退回到这道题的边界不变量。

\subsubsection{LeetCode 714 买卖股票含手续费}

\textbf{题面模型。} 手续费改变交易收益，仍可用持有和不持有状态。

\textbf{暴力瓶颈。} 卖出时扣费或买入时扣费都可以，但不要重复扣。

\textbf{边界反例。} 手续费为零、价格震荡、小利润不足手续费。

\textbf{变体迁移。} 这是状态机比局部贪心更稳的例子。

\subsubsection{LeetCode 739 每日温度}

\textbf{题面模型。} 每一天等待右侧第一个更高温度。

\textbf{暴力瓶颈。} 单调递减栈保存尚未找到答案的下标。

\textbf{边界反例。} 没有更高温、相等温度、递增、递减。

\textbf{变体迁移。} 下一个更大元素与它同型，只是输出值或距离不同。

\subsubsection{LeetCode 743 网络延迟时间}

\textbf{题面模型。} 有向正权图从起点到所有节点的最短路最大值。

\textbf{暴力瓶颈。} Dijkstra 求每个节点最短距离，最后取最大，若有无穷则不可达。

\textbf{边界反例。} 节点编号从一开始、不可达、重边、过期堆元素。

\textbf{变体迁移。} 边权相同才可用普通 BFS。

\subsubsection{LeetCode 752 打开转盘锁}

\textbf{题面模型。} 四位密码是状态，每次转动一位形成边。

\textbf{暴力瓶颈。} BFS 从起点扩展，死亡状态不能入队，访问状态避免重复。

\textbf{边界反例。} 起点死亡、目标就是起点、数字环绕、死亡集合很大。

\textbf{变体迁移。} 双向 BFS 能明显减少状态展开。

\subsubsection{LeetCode 763 划分字母区间}

\textbf{题面模型。} 每个片段必须覆盖其中所有字符的最后出现位置。

\textbf{暴力瓶颈。} 扫描维护当前片段最远右端，到达右端即可切分。

\textbf{边界反例。} 单字符、所有字符相同、字符多次跨段。

\textbf{变体迁移。} 这是最早合法切分贪心，能最大化片段数量。

\subsubsection{LeetCode 875 爱吃香蕉的珂珂}

\textbf{题面模型。} 速度越大，总耗时越小，答案具有单调可行性。

\textbf{暴力瓶颈。} 二分速度，检查函数用向上取整累加小时数。

\textbf{边界反例。} 速度下界、堆很大、h 等于堆数、乘法加法溢出。

\textbf{变体迁移。} 船运包裹和分割数组都是同类最小可行上限。

\subsubsection{LeetCode 973 最接近原点的 K 个点}

\textbf{题面模型。} 距离只需比较平方，避免开方误差和成本。

\textbf{暴力瓶颈。} 大小为 K 的大顶堆保存当前最近 K 个点中最远者。

\textbf{边界反例。} K 为一、K 为全部、距离相同、坐标平方溢出。

\textbf{变体迁移。} 快速选择可平均线性，但堆更适合数据流。

\subsubsection{LeetCode 994 腐烂的橘子}

\textbf{题面模型。} 所有初始腐烂橘子同时作为 BFS 源点。

\textbf{暴力瓶颈。} 按层扩散，分钟数等于最后一个新鲜橘子被首次访问的层数。

\textbf{边界反例。} 没有新鲜橘子、没有腐烂源、隔离新鲜橘子、空格阻断。

\textbf{变体迁移。} 多源 BFS 也适用于最近零距离、地图扩散等题。

\subsubsection{LeetCode 1011 在 D 天内送达包裹}

\textbf{题面模型。} 容量越大，需要天数越少，存在最小可行容量。

\textbf{暴力瓶颈。} 下界是最大包裹，上界是总重量，检查函数按顺序贪心装船。

\textbf{边界反例。} D 为一、D 等于包裹数、单个包裹超过猜测容量。

\textbf{变体迁移。} 它和分割数组最大值都是连续分段最大和模型。

\subsubsection{LeetCode 1143 最长公共子序列}

\textbf{题面模型。} 状态表示两个前缀的最长公共子序列长度。

\textbf{暴力瓶颈。} 末尾相同取左上加一，不同取上方和左方最大。

\textbf{边界反例。} 空串、完全相同、无公共字符、子序列不是子串。

\textbf{变体迁移。} 若要求还原序列，需要从表右下角反向追踪选择。

\subsubsection{LeetCode 1288 删除被覆盖区间}

\textbf{题面模型。} 被覆盖要求另一区间起点不晚且终点不早。

\textbf{暴力瓶颈。} 按起点升序、终点降序排序，扫描维护最大右端。

\textbf{边界反例。} 同起点区间、完全覆盖、无覆盖、端点相等。

\textbf{变体迁移。} 终点降序是关键，否则同起点小区间会先出现造成误判。

\subsubsection{LeetCode 1631 最小体力消耗路径}

\textbf{题面模型。} 路径代价是沿途最大高度差，不是边权和。

\textbf{暴力瓶颈。} Dijkstra 的松弛改成取当前代价和新边代价的最大值。

\textbf{边界反例。} 单格、平坦网格、必须绕路、多个相同代价。

\textbf{变体迁移。} 也可二分体力上限后 BFS 检查可达。

